% PROGRAMA DE PÓS-GRADUAÇÃO EM ECONOMIA APLICADA
% UNIVERSIDADE FEDERAL DO RIO GRANDE - FURG

% ===========================================================

\documentclass[a4paper, 11pt]{article}

%%%%%%%%%%%%%%%%%%%%%%%%% PACOTES %%%%%%%%%%%%%%%%%%%%%%%%%%%%%%%%%%

\usepackage[margin =0.5in]{geometry}
\usepackage[utf8]{inputenc}
\usepackage[T1]{fontenc}
\usepackage{graphicx}
\usepackage{xcolor} 
\usepackage{lipsum}	 
\usepackage{amsfonts,amsmath,amssymb, amsthm} 
\usepackage{hyperref}
\usepackage{layout}
% Configuração da Fonte
\usepackage{times} % Fonte: Times New Roman

% PACOTES ESSENCIAIS
\usepackage{caption}  

%\setlength{\footskip}{0pt} % Espacio entre texto y pie de página
\renewcommand{\footnoterule}{\vspace{5pt}\hrule width 0.3\linewidth\vspace{5pt}}
\setlength{\textheight}{770pt}
\setlength{\headsep}{10pt}
\setlength{\topmargin}{-60pt}

\begin{document} %Início do Documento
%\layout
%\newpage 
\- \vspace{-0.3in}
\par\noindent\includegraphics[width = \textwidth]{Imagens/Header.png}

%%%%%%%%%%%%%%%%%%%%%%%% CONFIG. DO TÍTULO %%%%%%%%%%%%%%%%%%%%%%%%
%%\vspace{n}

\begin{center} %Inicia a centralização
\textbf{Analisis Funcional | 2025-I}    \\ {\small
Nestor Heli Aponte Avila\(^1\) \\ 
\href{mailto:n267452@dac.unicamp.br}{\url{n267452@dac.unicamp.br}}} %\\
%Prof. Piper Pimienta \\
%\href{mailto:pipe@dac.unicamp.br}{\url{pipe@dac.unicamp.br}}}

\end{center} %Finaliza a centralização

%%%%%%%%%%%%%%%%%%%%%%%%% PREAMBULO %%%%%%%%%%%%%%%%%%%%%%%%%%%%%%%%%%

\theoremstyle{definition}
\newtheorem*{definition}{\rotatebox{45}{\large \(\square\)}}

\theoremstyle{plain}
\newtheorem*{lemma}{{\scriptsize \(\square\)}}
\newtheorem*{proposition}{{\large \(\square\)}}
\newtheorem*{theorem}{{\large \(\blacksquare\)}}

\makeatletter
% Eliminar la puntuación del estilo 'plain' y 'definition'
\def\@thm@headpunct{} % Elimina el punto de los títulos de los teoremas

% Modificar el estilo de 'plain' y 'definition'
\def\th@plain{%
  %\thm@headfont{\normalfont\bfseries} % Mantener la fuente normal, sin negrita
  %\thm@notefont{\normalfont} % El texto de la nota (si la hay) no será en negrita
  \thm@headpunct{} % Sin puntuación adicional
}

\def\th@definition{%
  %\thm@headfont{\normalfont\bfseries} % Fuente normal para el título
  %\thm@notefont{\normalfont} % El texto de la nota (si la hay) no será en negrita
  \thm@headpunct{} % Sin puntuación adicional
}

\def\th@remark{%
  %\thm@headfont{\normalfont\bfseries} % Fuente normal para el título
  %\thm@notefont{\normalfont} % El texto de la nota (si la hay) no será en negrita
  \thm@headpunct{} % Sin puntuación adicional
}
\makeatother

\theoremstyle{remark}
\newtheorem*{example}{Ejemplo}
\newtheorem*{exercise}{Ejercicio}
\newtheorem*{note}{\(*\)}

\newcommand{\R}{\mathbb{R}}
\newcommand{\N}{\mathbb{N}}
\newcommand{\C}{\mathbb{C}}
\newcommand{\K}{\mathbb{K}}
\newcommand{\Li}{\mathcal{L}}
\newcommand{\E}{\vspace{0.2in}}
\newcommand{\ab}{{\rotatebox{0}{\scriptsize$\lor$}}}
\newcommand{\ce}{{\text{\ \large \rotatebox{90}{$\triangleleft$}}}}

%%%%%%%%%%%%%%%%%%%%%%%%%% REFERÊNCIAS %%%%%%%%%%%%%%%%%%%%%%%%%%%%%%%%

\section*{Espacios Vectoriales Normados}

\subsection*{Definición y Ejemplos}

\begin{definition}
    Sea \(E\) un espacio vectorial sobre \(\K\). Una función \(\|\cdot\|: E\rightarrow \R^+ \) es \emph{norma} sii \(\forall x,y\in E \ \text{y}\ \alpha\in \K\),   
    \begin{itemize}
        \item \(\|x\|\geq 0\) e \(\|x\|=0 \Leftrightarrow x = 0\). 
        \item \(\|\alpha x\| = |\alpha|\|x\|\). 
        \item \(\|x+y\|\leq \|x\|+\|y\|\). 
    \end{itemize} 
\end{definition}

\begin{note}
    \((E,\|\cdot\|)\) espacio normado implica espacio métrico, vale la teoría existente para ellos en particular convergencia. 
\end{note}

\begin{definition}
    Sea \((x_n)\subset E\). Decimos que \(x_n\to x \in E\) sii \(\lim\limits_{n\to\infty} \|x_n - x\| = 0\).  
\end{definition}

\E

\hrule 
\begin{exercise}
    Las operaciones algebraícas en \(E\) son funciones continuas. 
\end{exercise}
\hrule 

\E

\begin{definition}
    \((E,\|\cdot\|)\) es \emph{Banach} sii es completo con la métrica inducida por la norma. 
\end{definition}

\E  

\hrule 
\begin{example}
    \((\R,\|\cdot\|_2)\) y \((\C,\|\cdot\|_2)\) son espacios de Banach. 
\end{example}
\hrule 

\E  

\begin{proposition}
    Sea \(E\) Banach y \(F\leq E\) un subespacio vectorial, entonces \(F\) es Banach sii \(F\) es cerrado en \(E\). 
\end{proposition}

\E 

\hrule 
\begin{example}
    Sea \(B(X):= \{f:X\rightarrow \K \ |\ f \ \text{es acotada}\}\) e \(\|f\|_\infty := \sup |f(x)|\), entonces \(B(X)\) es Banach. Sea \([a,b]\subset \R\), conjunto compacto, observe que \(C^0[a,b]\leq B[a,b]\) normado.   
\end{example}
\begin{exercise}
    Complete los detalles del ejemplo, y muestre que \(C^0[a,b]\) es Banach.
\end{exercise}
\begin{example}
    Considere \(C^1[a,b] \leq C^0[a,b]\), no es Banach. Sin embargo con \(\|f^{(1)}\|_{\infty^1}:= \|f\|_\infty + \|f^{(1)}\|_\infty\) si que lo es. En general \(C^k[a,b]\) es Banach con \(\|f^{(k)}\|_{\infty^k} = \sum_{i=0}^k \|f^{(i)}\|_\infty\). 
\end{example}
\hrule 

\E

\begin{proposition}
    Si \(B = \{x_1,\ldots,x_n\}\) es una base l.i. de \(E\), entonces \(\exists c>0, \forall a \in \K^n\) tal que \(\left\|\sum a_ix_i\right\|\geq c \sum |a_i|\). 
\end{proposition}
\begin{theorem}
    Todo espacio \(E\) tal que \(\dim{(E)}<\infty\) es Banach. Consecuentemente, también lo son todos los \(F\leq E\) cerrados. 
\end{theorem}

\E

\hrule
\begin{example}
    Sea \(c_0=\{(a_k) \subset \K : a_k \rightarrow 0\}\) con las operaciones usuales y \(\|(a_k)\|_\infty := \sup |a_k|\). Así dado \(c_0\) es Banach. 
\end{example}
\begin{example}
    Sea \(c_{00}:= \{(a_k)\in c_0 : \exists n_0\in \N \text{ tal que } a_k = 0\text{ para } k\geq n_0\}\). No cerrado, no Banach. 
\end{example}
\hrule 

\subsection*{Espacios \(L_p,\ \ell_p\)}

\begin{definition}
    Sea \(\Li(X,\Sigma,\mu)\) espacio de medida. Consideremos \(\Li_p(X,\Sigma,\mu)\) el subespacio de las \(f:X\rightarrow\K\) tales que para \(1\leq p<\infty \) el valor \(\|f\|_p := \displaystyle \left(\int_X |f|^p\ d\mu\right)^{\frac{1}{p}}< \infty \). 
\end{definition}
\begin{theorem}
    \emph{Hölder.} Sean \(p,q>1\) tales que \(1/p + 1/q = 1\). Si \(f\in \Li_p\) e \(g\in \Li_q\) entonces \(fg\in \Li_1\) y \(\|fg\|_1\leq \|f\|_p\cdot\|g\|_q\).  
\end{theorem}
\begin{theorem}
    \emph{Minkowski.} Si \(f,g\in \Li_p\) entonces \(f+g\in \Li_p\) y \(\|f+g\|_p\leq \|f\|_p+\|g\|_p\) 
\end{theorem}
\begin{note}
    Note que \(\|f\|_p = 0 \nRightarrow f=0\). Ello motiva la siguiente consideración.  
\end{note}
\begin{definition}
    Sean \(f,g\in \Li_p\). Decimos que \(f\sim g\) si \(f=g\) \(\mu\)-casi siempre, es decir, \(\exists A\in \Sigma \) tal que \(\mu(A)=0\) y \(f(x)=g(x)\) para todo \(x\in X\setminus A\). 
\end{definition}
\begin{theorem}
    El conjunto \(L_p = \Li_p/\hspace{-0.1cm}\sim\) con las operaciones \([f]+[g] = [f+g]\), \([cf]=c[f]\) y norma \(\|[f]\|_p := \|f\|_p\) es Banach. 
\end{theorem}
\begin{definition}
    Sea \(\Li_\infty\) el espacio de las \(f\) medibles acotadas \(\mu\)-cuasi siempre\footnote{\(|f(x)|\leq k < \infty\) para cada \(x\in X\setminus N\) donde \(\mu(N\in \Sigma)=0\).}. Para cada \(x\in X\setminus N\) sean \(S_f(N) = \sup|f(x)|\) y \(\|f\|_\infty:= \inf S_f(N)\) de los \(N\in \Sigma\) tales que \(\mu(N) = 0\). 
\end{definition}
\begin{theorem}
    El espacio \(L_\infty = \Li_\infty/\hspace{-0.1cm}\sim \) con \(\|[f]\|_\infty := \|f\|_\infty\) es Banach. 
\end{theorem}

\E

\hrule 
\begin{example}
    Sean \(p\geq 1\) y \(\ell_p = \left\{(a_k)\subset \K : \displaystyle \sum |a_k|^p < \infty\right\}\). Considere \(\Sigma = \wp(\N)\) y \(\mu_c\) la medida de conteo. El espacio \(\ell_p\) coincide con \(L_p(\N,\wp(\N), \mu_c)\), en este caso las operaciones son simplemente las usuales de sucesiones y \(\displaystyle \|(a_n)\|_p = \left(\sum |a_k|^p\right)^{\frac{1}{p}}\). Entonces \(\ell_p\) es Banach.
\end{example}
\hrule 

\E

\begin{proposition}
   \emph{Hölder-Minkowski en sucesiones.} Para \(p,q>1\) tales que \(1/p+1/q = 1\) y \(\forall n\in \N\) vale que 
   \[\sum^n |a_ib_i|\leq \left(\sum^{n}|a_i|^p\right)^{\frac{1}{p}} \cdot \left(\sum^{n}|b_i|^q\right)^{\frac{1}{q}}\text{\ \ y \ \ } \left(\sum^{n}|a_i + b_i|^p\right)^{\frac{1}{p}}\leq \left(\sum^{n}|a_i|^{p}\right)^{\frac{1}{p}} + \left(\sum^{n}|b_i|^{p}\right)^{\frac{1}{p}}\]
\end{proposition}

\E

\hrule
\begin{example}
    El espacio \(\ell_\infty = \{(a_k)\subset \K: \sup |a_k| <\infty\} \) con la norma \(\|(a_k)\|_\infty = \sup |a_k|\) es Banach. Esto es directo observando que \(\ell_\infty = L_\infty(\N,\wp(\N),\mu_c) = B(\N)\)\footnote{Verificando alguna de las dos igualdades.}.  
\end{example}
\hrule 

\subsection*{Compacidad}

\begin{definition}
    \(A\subseteq X\) es compacto sii todo cubrimiento abierto de \(A\) admite un subcubrimiento finito. 
\end{definition}
\begin{note}
    En espacios métricos vale decir que toda \((a_k)\subset A\) admite una subsucesión \((a_{k_i})\) tal que \(a_{k_i} \rightarrow a \in A\). 
\end{note}
\begin{proposition}
    Sea \((E,\|\cdot\|)\) con \(\dim(E)<\infty\), entonces los compactos de \(E\) son precisamente los cerrados y acotados. 
\end{proposition}
\begin{note}
    \emph{Colorario.} La \emph{bola unitaria} \(B_E := \{x\in E : \|x\|\leq 1 \}\) es compacta en espacios de dimensión finita. 
\end{note}
\begin{proposition}
    \emph{Riesz.} Si \(M<E\) cerrado y \(\theta \in (0,1)\), entonces \(\exists y\in E\setminus M,\forall x\in M\) tal que \(\|y\|=1\) y \(\|y-x\|\leq \theta \). 
\end{proposition}
\begin{theorem}
    \(\dim(E)<\infty\) sii \(B_E\) es compacta en \(E\). 
\end{theorem}

\vfill 

\subsection*{Espacios Separables}

\begin{definition}
    \(E\) es separable sii \(\exists D\subset E\) enumerable y denso en \(E\). 
\end{definition}

\E

\hrule
\begin{example}
    Espacios con \(\dim (E)<\infty\) son separables. Caso de \(\mathbb{Q}\subset \R\). 
\end{example}
\hrule 

\E

\begin{proposition}
    \(E\) es separable sii \(\exists A\subset E\) enumerable tal que \(\langle A\rangle\) es denso en \(E\). 
\end{proposition}

\E

\hrule 
\begin{example}
    \(c_0\) y \(\ell_p \) son separables, mientras \(\ell_\infty\) no lo es. 
\end{example}
\hrule

\E

\begin{theorem}
    \emph{Aproximación de Weierstrass.} \(f:[a,b]\rightarrow \K\) continua \(\Rightarrow \forall\epsilon>0,\forall x\in [a,b], \exists P:\K\rightarrow \K \) tal que \(|P(x) - f(x)|<\epsilon \). 
\end{theorem}

\E

\hrule 
\begin{example}
    \(C[a,b]\) es separable. Hint:\(\langle t \rangle \). 
\end{example}
\begin{example}
    \(L_p[a,b]\) es separable. Hint: continuas y polinomios. 
\end{example}
\hrule 

\E

\begin{proposition}
    Si \(E\) es separable entonces \(F\leq E\) también lo es. 
\end{proposition}

\section*{Operadores Lineales}

\begin{definition}
    Un \emph{operador lineal continuo} es una función \(T:E\to F \) que verifica lo siguiente, 
    \begin{itemize}
        \item \(\forall x,y \in E,\forall \alpha \in \K \ \text{tenemos } T(\alpha x+y) =  \alpha T(x) + T(y)\). 
        \item \(\forall x_0 \in E,\forall \epsilon > 0, \exists \delta > 0\) tal que \(\|x-x_0\|<\delta \Rightarrow \|T(x)-T(x_0)\|<\epsilon\). 
    \end{itemize}
 \end{definition}
\newcommand{\Lc}{\mathcal{L}}
\begin{note}
    \(\Lc(E,F)=\{T:E\to F\ |\ T \text{ es lineal continuo}\}\) es espacio vectorial sobre \(\K\). Si \(F = \K\) entonces \(\Lc(E,\K)=E'\), el espacio dual, cuyos elementos son funciones.   
\end{note}
\begin{definition}
    \(E\cong F\) sii \(\exists T\in \Lc(E,F)\) biyectivo cuyo inverso \(T^{-1}\in \Lc(F,E)\). 
\end{definition}
\begin{definition}
    Una función \(f:E\to F \) es una \emph{isometría} sii \(\forall x\in E\) tenemos \(\|f(x)\| = \|x\|\).      
\end{definition}
\begin{note}
    Si \(f\) es lineal entonces es una \emph{isometría lineal}. Si \(f\) es un isomorfismo entonces le llamamos \emph{isomorfimo isométrico}. 
\end{note}

\E 

\hrule 
\begin{exercise}
    Toda isometría lineal es inyectiva y continua. 
\end{exercise}
\hrule 

\subsection*{Caracterización}

\begin{definition}
    Una función \(f:M\to N \) es \emph{Lipschitz} si \(\exists L>0, \forall x,y\in M\) tal que \(\|f(x)- f(y)\|\leq L \|x-y\|\).  
\end{definition}
\begin{definition}
    \(f\) es \emph{uniforme continua} si \(\forall \epsilon >0,\exists \delta >0,\forall x,y\in M\) tal que \(\|x-y\|<\delta \Rightarrow \|f(x)-f(y)\|<\epsilon\). 
\end{definition}
\begin{note}
    Lipschitz \(\Rightarrow\) uniforme continua \(\Rightarrow\) continua \(\Rightarrow\) continua en \(x_0\). 
\end{note}
\begin{theorem}
    Sea \(T\in \Lc(E,F)\). \(T\) es Lipschitz \(\Leftrightarrow\) \(T\) es uniforme continuo \(\Leftrightarrow\ T\) es continuo \(\Leftrightarrow\ \exists x_0\in E\) tal que \(T\) es continuo en \(x_0\) \(\Leftrightarrow\ T\) es continuo en \(0\)  \(\Leftrightarrow\ \sup\{\|T(x)\| : x\in B_E\}<\infty\) \(\Leftrightarrow \) \(\exists C\geq 0, \forall x\in E\) tal que \(\|T(x)\|\leq C \|x\|\). 
\end{theorem}
\begin{note}
    \emph{Colorario.} \(T\in \Lc(E,F)\) biyectivo es isomorfismo sii \(\exists C_1,C_2>0, \forall x\in E\) tal que \(C_1\|x\|\leq \|T(x)\| \leq C_2 \|x\|\). 
\end{note}
\begin{proposition}
    Sean \(E,F\) espacios normados, entonces \((a) \ \|T\| = \sup\{\|T(x)\|: x\in B_E\}\) es norma en \(\Lc(E,F)\); \((b) \ \forall T\in \Lc(E,F),\forall x\in E \) tenemos \( \|T(x)\| = \|T\|\cdot\|x\|\) y \((c) \ F\text{ Banach } \Rightarrow \Lc(E,F) \text{ Banach}\). 
\end{proposition}
\begin{note}
    \emph{Colorario (c).} \(E'\) es Banach. 
\end{note}

\E 

\hrule 
\begin{example}
    El operador identidad \(1_E: E \to E \in \Lc(E,E)\) para el cual \(\|1_E\|=1\); Operador nulo \(O:E\to F \in \Lc(E,F)\) que envía \(x\mapsto 0_F\) tenemos \(\|O\| = 0\). 
\end{example}
\begin{example}
    Sea \(\varphi \in E'\) e \(y\in F\). Sea \(\varphi \otimes y: x \mapsto \varphi(x)y \in \Lc(E,F)\) y tiene norma \(\|\varphi\|\cdot\|y\|\).  
\end{example}
\begin{example}
    Sea \((b_n) \in \ell_p\). Considere \(T:\ell_\infty \to \ell_p\) tal que \((a_n)\mapsto (a_nb_n)\) \footnote{\(T\) es llamado \emph{operador diagonal} por \((b_n)\).}. 
\end{example}
\begin{example}
    Sea \(g\in L_p[0,1]\). Como en el ejemplo anterior considere \(T:C[0,1] \to L_p[0,1], \ T(f) = fg\). El operador \(T\in \Lc(C[0,1],L_p[0,1])\) y es llamado \emph{operador multiplicación}. 
\end{example}
\begin{exercise}
    \(T\) lineal en \(E\) con \(\dim(E)<\infty \Rightarrow \ T \) continuo. En dimensión infinita no siempre es cierto. 
\end{exercise}
\begin{example}
    Sea \(\mathcal{P}[0,1]\subset C[0,1]\) con la norma \(\|\cdot\|_\infty \). El operador derivación es lineal, suponga continuo, entonces \(\exists C, \forall p \in \mathcal{P}[0,1] \) tal que \(\|T(p)\|_\infty \leq C\|p\|_\infty\). Sea \(f_n = t^n\), tenemos \(n = \|f_n^{'}\|_\infty = \| T(f_n)\|_\infty \leq C\|f_n\|_\infty = C \). 
\end{example}
\hrule 

\subsection*{Teorema Banach-Steinhaus}

\begin{theorem}
    \emph{Baire.} Sea \(M\) espacio métrico completo y \(\left(F^\ce_n\right) \subseteq M\) tal que \(M=\bigcup F_n\). Entonces \(\exists n_0\in \N\) tal que \(\overset{\circ}{F_{n_0}} \neq \emptyset\).    
\end{theorem}
\begin{theorem}
    \emph{Banach-Steinhaus.} Sean \(E\) Banach, \(F\) espacio normado y \((T_i)\) una sucesión de operadores en \(\Li(E,F)\) tales que \(\forall x\in E,\exists C_x <\infty\) tal que \(\sup \|T_i(x)\| < C_x\). Entonces \(\sup \|T_i\| < \infty\).  
\end{theorem}
\begin{note}
    \emph{Colorario.} Sea \((T_n)\subset \Li(E,F)\). Si \(\forall x\in E\) la sucesión \((T_n(x)) \to y\in F\) entonces \(T(x) = \lim T_n(x) \in \Li(E,F)\).  
\end{note}

\E 

\hrule
\begin{example}
    \((x,y)\mapsto \frac{xy}{x^2+y^2}\), \((0,0) \mapsto 0\) es continua en  \(\R^2\setminus \{0\}\), en aplicaciones \emph{bilineales} no existen cosas así. 
\end{example}
\hrule 

\E

\begin{definition}
    Sean \(E_1, E_2 \text{ y } F \) espacios vectoriales. Una aplicación \(B:E_1\times E_2 \to F\) es \emph{bilineal} sii \(\forall x_1\in E_1,\forall x_2\in E_2 \text{ fijos, los operadores }B(x_1,\cdot): E_2 \to F\) y \(B(\cdot,x_2): E_1\to F\) son lineales. 
\end{definition}

\begin{note}
    \emph{Colorario.} Si \(E_2\) es completo y \(B:E_1\times E_2 \to F \) es bilineal y continuo a trozos entonces \(B\in \Li(E_1\times E_2, F)\).    
\end{note}

\E

\hrule 
\begin{example}
    \(\forall n\in\N\) sea \(\varphi_n: c_{00} \ni (a_j) \mapsto na_n \in \K\). Es claro que \((\varphi_n)\subset (c_{00})'\) y \(\|\varphi_n\|=n\), aquí \(\forall x\in c_{00} \text{ se tiene }\sup \|\varphi_n(x)\|<\infty\), sin embargo, \(\sup\|\varphi_n\|=\infty\). 
\end{example}
\hrule 

\E

\subsection*{Teorema de la Aplicación Abierta}

\begin{definition}
    Nos referimos a \(B_E(x_0;r)= \{x\in E: \|x-x_0\|<r\}\) como bola abierta en \(E\) centrada en \(x_0\) de radio \(r>0\). 
\end{definition}
\begin{proposition}
    Sean \(E\) Banach, \(F\) espacio normado y \(F\leftarrow E : T \in \Li(E,F)\). Si existieran \(R,r >0\) tales que \(\overline{T(B_E(0;R))}\supseteq B_F(0;r)\) entonces \(T(B_E(0;R)) \supseteq B_F\left(0;\frac{r}{2}\right)\).  
\end{proposition}
\begin{theorem}
    \emph{Aplicación Abierta.} Sean \(E\text{ e }F\) Banach. Si \(F \leftarrow E:\overset{\twoheadrightarrow}{T}\in \Li(E,F)\) entonces \(T\) es una aplicación abierta.  
\end{theorem}
\begin{note}
    \emph{Colorario.} En particular si \(T\) es una biyección entonces \(E\cong F\).  
\end{note}

\E

\hrule 
\begin{exercise}
    Muestre que  \(T: c_{00}\to c_{00}\) tal que \((a_n)\mapsto \left(\frac{a_n}{n}\right)\) es lineal, continuo y biyectivo. 
\end{exercise}
\begin{example}
    En el ejercicio anterior \(T^{-1}\) no es continuo.  
\end{example}
\begin{example}
    Todo subespacio \(F^\ce\leq C[0,1]\) tal que \(\dim (F) = \infty \) tiene al menos una función \(f\) tal que \(f\notin C^1[0,1]\). Hint: Contradicción -- Aplicación Abierta -- Teorema de Riesz. 
\end{example}
\hrule 

\E

\begin{definition}
    Sean \(E \text{ e } F\) espacios normados y \(T:E\to F\) lineal. El \emph{gráfico} de \(T\) es el conjunto, 
    \[G(T) = \{(x,T(x)):x\in E\}\subseteq E\times F. \]
\end{definition}
\begin{theorem}
    \emph{Gráfico Cerrado.} Sean \(E \text{ e } F\) Banach y \(T:E\to F\) lineal. El operador \(T\) es continuo sii \(G(T)\) es cerrado en \(E\times F\). 
\end{theorem}

\E

\hrule 
\begin{exercise}
    Si \(T\) no es continuo una de las implicaciones en el Teorema del Gráfico Cerrado continua valiendo. 
\end{exercise}
\begin{example}
    Sean \(E\) Banach y \(T:E\to E'\) lineal \emph{símetrico}, es decir, \(\forall x,y\in E \text{ tenemos }T(x)(y)=T(y)(x)\). El operador \(T\) es continuo. Hint: Gráfico Cerrado. 
\end{example}
\hrule 

\section*{Teoremas de Hahn-Banach}

\begin{proposition}
    \emph{Lemma de Zorn.} Todo conjunto parcialmente ordenado, no vacío y en el cual todo subconjunto totalmente ordenado tiene cota superior, tiene un elemento máximal. 
\end{proposition}

\E

\hrule
\begin{exercise}
    Lemma de Zorn \(\Leftrightarrow \) Axioma de Elección. 
\end{exercise}
\hrule 

\E

%\begin{theorem}
%    \emph{Hanh-Banach (en \(\R\)).} Sean \(E\) un espacio (sobre \(\R\)) normado y \(p:E\to \R\) una función tal que, 
%    \begin{itemize}
%        \item \(\forall a>0,\forall x\in E\) se tiene \(p(ax) = ap(x)\). 
%        \item \(\forall x,y\in E\) su cumple \(p(x+y) \leq p(x)+p(y)\). 
%    \end{itemize} 
%    Sean tambien \(G\leq E \) y \(\varphi : G \to \R\) un operador lineal tal que \(\forall x\in G \text{ se tiene }\varphi(x)\leq p(x)\). Entonces \(\exists \overset{\sim }{\varphi} : E \to \R\) lineal que extiende a \(\varphi\), es decir, \(\overset{\sim }{\varphi}(x)\Big|_{G} = \varphi(x) \) y que además satisface \(\forall x\in E \text{ que }\overset{\sim}{\varphi}(x)\leq p(x)\). 
%\end{theorem}
\begin{theorem}
    \emph{Hanh-Banach (en \(\K\)).} Sean \(E\) un espacio (sobre \(\K\)) normado y \(p:E\to \R\) una función tal que, 
    \begin{itemize}
        \item \(\forall a>0,\forall x\in E\) se tiene \(p(ax) = |a|p(x)\). 
        \item \(\forall x,y\in E\) su cumple \(p(x+y) \leq p(x)+p(y)\). 
    \end{itemize} 
    Si \(G\leq E \) y \(\varphi : G \to \K\) es un operador lineal tal que \(\forall x\in G \text{ se tiene }|\varphi(x)|\leq p(x)\), entonces \(\exists \overset{\sim }{\varphi} : E \to \K\) lineal que extiende a \(\varphi\), es decir, \(\overset{\sim }{\varphi}(x)\Big|_{G} = \varphi(x) \) y que además satisface \(\forall x\in E \text{ que }|\overset{\sim}{\varphi}(x)|\leq p(x)\). 
\end{theorem}
\begin{note}
    \emph{Colorario(s).} \begin{itemize}
        \item Si \(\varphi \) es continuo entonces \(\overset{\sim}{\varphi}\) también y \(\|\varphi\| = \|\overset{\sim}{\varphi}\|\). 
        \item Si \(E\) es un espacio normado entonces \(\forall x_0\in E \setminus \{0\},\exists \varphi\in E'\) tal que \(\|\varphi\| = 1\) y \(\varphi(x_0) = \|x_0\|\). 
        \item Si \(E\neq \{0\}\) y \(x\in E \) entonces \(\|x\| = \sup\{ |\varphi(x)|: \varphi \in E' \text{ y } \|x\|\in B_E\} \) cuyo valor alcanza.   
    \end{itemize}
\end{note}

\subsection*{Versiones Vectoriales del Teorema de Hahn-Banach}

\begin{definition}
    Sea \(E\) Banach y \(P\in \Li(E,E)\) es una \emph{proyección} sii \(P^2 = P \circ P = P \). 
\end{definition}
\begin{note}
    Si \(P\neq 0\) entonces \(\|P\|\geq 1 \). 
\end{note}
\begin{proposition}
    Sea \(F\leq E\). Entonces son equivalentes \((a)\ \exists P\in \Li(E,E)\) proyección tal que \(P(E) = F   \) e \((b) \ F^\ce \leq E \) e \(\exists G^\ce \leq E  \) tal que \(E = F\oplus G\). 
\end{proposition}
\begin{note}
    \(F = \{x\in E : P(x) = x\}\) e \(G = \ker(P)\). 
\end{note}
\begin{definition}
    \(F\leq E \) es complementado si satisface alguna \((a)\) o \((b)\). Es \(\lambda\)-complementado si \(\|P\| = \lambda\). 
\end{definition}

\E

\hrule 
\begin{example}
    Todo \(F\leq E\) con \(\dim (F)<\infty \) es complementado. Hint: Base + $a_j$ + Hahn-Banach.   
\end{example}
\begin{example}
    Siendo \(E\) y \(F\) Banach. La proyección \(E\times F \ni (x,y)\mapsto (x,0) \in E\times \{0\}\) deja ver que \(E\) es $1$-complementado. 
\end{example}
\begin{example}
    SI existen espacios cerrados que son no complementados, créditos a Murray \((\ell_p)\) y Phillips \((\ell_\infty)\). 
\end{example}
\hrule 

\E

\begin{proposition}
    Sean \(G\) Banach, \(F\leq E \) complementado e \(T \in \Li(F,G)\), entonces \(\exists \overset{\sim}{T} \in \Li(E,G)\). 
\end{proposition}
\begin{proposition}
    Si \(F\leq E\) no complementado, entonces \(\not\exists T\in \Li(E,F), \forall x\in F \) tal que \(T(x) = x \).  
\end{proposition}
\begin{note}
    La identidad no puede ser extendida continuamente a \(E\). 
\end{note}
\begin{theorem}
    \emph{Phillips.} Sean \(F\leq E\) e \(T\in \Li(E,\ell_\infty)\). Entonces \(\exists \overset{\sim}{T}\in \Li(E, \ell_\infty)\), con \(\|\overset{\sim}{T}\| = \|T\|\). 
\end{theorem}
\begin{note}
   \emph{Colorario.}  Si \((\ell_\infty)^\ce \leq E\) entonces \(\ell_\infty\) es $1$-complementado en \(E\). 
\end{note}

\subsection*{Aplicaciones de Hahn-Banach a Espacios Separables}

\begin{proposition}
    Sean \(M^\ce \leq E\), \(y_0\in  E\setminus M\) y \(d = \textnormal{dist}(y_o,M)\). Entonces \(\exists \varphi \in E', \forall x\in M\) tal que \( \|\varphi\|=1 \), \(\varphi(y_0) = d\) y \(\varphi(x)=0\). 
\end{proposition}
\begin{theorem}
    Si \(E'\) es separable, entonces \(E\) también. 
\end{theorem}
\begin{proposition}
    \(\forall E\) separable se tiene \(E\cong F\leq \ell_\infty\) isómetricamente.  
\end{proposition}

\subsection*{Formas geométricas del Teorema de Hahn-Banach (\(\K = \R\)) }

\begin{definition}
    Sea \((V,+,\cdot)\neq \{0\}\). El subespacio \(W< V \) es \emph{hiperplano} sii \(W < W_1< V\) implica \(W_1= V\). 
\end{definition}
\begin{proposition}
    \(W<V\) es hiperplano sii \(\exists \varphi \neq 0: V \to \R \) lineal tal que \(W = \ker(\varphi)\). 
\end{proposition}
\begin{note}
    Si \(H<V\) es un hiperplano entonces \(v_0 + H = \{v\in V: \varphi(v) = \varphi(v_0) = a\in \R\}\) es un \emph{hiperplano afín}. 
\end{note}
\begin{proposition}
    \(H^\ce <V\) sii \(\varphi \) es continua. 
\end{proposition}
\begin{definition}
    Sea \(C^\ab \ni 0\subseteq E\) convexo. Llamamos \emph{funcional de Minkowski} a la aplicación \(p_C:E\to \R\) que envía \(x\mapsto \inf \left\{a>0: \frac{x}{a} \in C \right\}\). 
\end{definition}
\begin{proposition}
    El funcional de Minkowski verifica \(\forall b>0, \forall x,y\in E\) que:  \((a)\ p_C(bx) = bp_C(x)\); \((b) \ C = \{x\in E : p_C(x)<1\}\); \((c) \ \exists M>0 \text{ tal que } 0\leq p_C(x)\leq M \|x\|\) e \((d)\ p_C(x+y) \leq p_C(x)+p_C(y)\).  
\end{proposition}
\begin{proposition}
    Sean \(\emptyset \neq C^\ab \subset E \) conexo e \(x_0\in E\setminus C\), entonces \(\exists \varphi \in E', \forall x\in C\) tal que \(\varphi(x)<\varphi(x_0)\).
\end{proposition}
\begin{proposition}
    Sean \(0\neq \varphi \in E'\) y \(\emptyset \neq A^\ab  \subset E\) convexo, entonces \(\varphi(A)^\ab\subseteq \R\).  
\end{proposition}
\begin{theorem}
    \emph{Primera Forma Geométrica del Teorema de Hahn Banach.} Sean \(\emptyset \neq A,B \subset E\) disjuntos. Si \(A^\ab\subset E\) entonces \(\exists \varphi \in E', \exists a \in \R, \forall x\in A, \forall y\in B\) se tiene \(\varphi(x) < a \leq \varphi(y)\). 
\end{theorem}
\begin{note}
    En este caso decimos que el hiperplano \([\varphi =  a]\) separa a \(A\) de \(B\). 
\end{note}
\begin{theorem}
    \emph{Segunda Forma Geométrica del Teorema de Hahn Banach.} Sean \(\emptyset \neq A,B \subset E\) disjuntos. Si \(A^\ce\subset E\) y \(B\) es compacto entonces \(\exists \varphi \in E', \exists a,b \in \R, \forall x\in A, \forall y\in B\) se tiene \(\varphi(x) < a <b \leq \varphi(y)\). 
\end{theorem}
\begin{note}
    Ahora decimos que \(\forall c\in (a,b)\) el hiperplano \([\varphi = c]\) separa estrictamente a \(A\) de \(B\). 
\end{note}
\begin{note}
    \emph{Colorario.} Si \(M^\ce\leq E\) entonces \(\forall x_0\in E\setminus M,\exists \varphi \in E', \forall x\in M\) se tiene \(\varphi(x_0)=1\) y \(\varphi(x) = 0\). 
\end{note}
\begin{note}
    \emph{Colorario.} Sea \(x_0\in M\leq E\), entonces \(x_0\in \overline{M}\) sii \(\forall \varphi \in E', \forall x\in M\) se tiene \(\varphi(x_0) = 0\) siempre que \(\varphi(x)=0\). 
\end{note}

%%%%%%%%%%%%%%%%%%%%%%%%%% REFERÊNCIAS %%%%%%%%%%%%%%%%%%%%%%%%%%%%%%%%

\begin{thebibliography}{9}

\bibitem{martin1966complex}
Martin, D. y Ahlfors, L.V. (1966). \textit{Complex Analysis}. New York: McGraw-Hill.

\end{thebibliography}

%%%%%%%%%%%%%%%%%%%%%%%%%% APÊNDICES %%%%%%%%%%%%%%%%%%%%%%%%%%%%%%%%

\end{document}
